\documentclass[12pt,a4paper]{article}

\usepackage[T1]{fontenc}
\usepackage[utf8]{inputenc}
\usepackage[catalan]{babel}
\usepackage[a4paper,margin=2.6cm]{geometry}
\usepackage{amsmath,amssymb}
\usepackage{hyperref}

\usepackage{newpxtext}
\usepackage{newpxmath}

\usepackage{setspace}
\usepackage{parskip}

\hypersetup{
  pdftitle={Guió de defensa},
  pdfauthor={Mario Vilar},
  hidelinks
}

\newcommand{\diapo}[2]{%
  \section*{Diapositiva #1: #2}%
  \addcontentsline{toc}{section}{Diapositiva #1: #2}%
}
\newcommand{\perque}[1]{%
  \par\smallskip\noindent\textbf{Per què serveix.}\ #1\par\smallskip%
}

\newenvironment{nota}{%
  \par\smallskip\noindent\textbf{Nota de ritme.}\ }{\par\smallskip}

\title{Guió de defensa}
\author{Mario Vilar}
\date{}

\begin{document}
\maketitle

\begin{nota}
Objectiu de ritme: màxim 20 minuts. Amb 29 diapositives principals, la majoria
han de durar entre 25 i 40 segons.
\end{nota}

\diapo{1}{Portada}
\perque{Situa el tema i la pregunta global sense gastar encara temps tècnic.}

Bon dia, em dic Mario Vilar i avui presentaré el meu TFG,
\emph{Models causals estructurals i resultats potencials}.

\diapo{2}{Índex}
\perque{Dona el mapa de la defensa i prepara el tribunal per al ritme ràpid.}

Primer motivaré el problema de la causalitat; després introduiré models causals
estructurals i identificació; a continuació passaré als resultats potencials; i
acabaré amb el cas de Berkeley i unes conclusions.

\diapo{3}{Què és la causalitat}
\perque{Motiva per què causalitat no és només veure regularitats.}

Aquí una fotografia d'un servidor amb l'estàtua de Hume a Edimburg. Hume és un
dels grans noms de la causalitat i va arribar a dir que la causalitat és
``el ciment del món''. Avui no faré història de la filosofia: definiré la
causalitat per oposició al que no és.

\diapo{4}{Què no és la causalitat}
\perque{Fixa la intuïció bàsica: una correlació pot ser completament espúria.}

Galton i Pearson creien que la causalitat era un cas molt concret de correlació.
Avui sabem que això no és suficient. Aquí tenim un exemple volgudament absurd:
la pluja a San Francisco i les rotatives a Rhode Island tenen una correlació
molt alta, però ningú defensaria una relació causal.

\diapo{5}{La jerarquia causal de Pearl}
\perque{Ordena les preguntes: observar, intervenir i imaginar no són el mateix.}

Pearl organitza les preguntes causals en tres nivells. El primer és
l'associació: què passa si observo \(X\)? El segon és la intervenció: què passa
si faig \(X\)? El tercer és el contrafactual: què hauria passat si...?

Com més pugem, més estructura causal i més hipòtesis necessitem. Aquesta és una
idea que tornarà al cas Berkeley.

\diapo{6}{Regla de la cadena}
\perque{Introdueix la factorització local que fa útil el graf.}

Aquí fixo una mínima notació. Una xarxa bayesiana és un graf i una distribució
que factoritza sobre aquest graf. La idea important és que el graf ens diu quina
informació és rellevant quan condicionem: els pares gràfics.

\diapo{7}{d-separació i camins}
\perque{Permet llegir independències condicionals directament del graf.}

La \(d\)-separació és la regla gràfica que permet llegir independències
condicionals. Cadenes i bifurcacions es bloquegen condicionant en el node
intermedi. Els col·lisionadors funcionen al revés: normalment bloquegen el camí,
però condicionar-hi el pot obrir.

\diapo{8}{Models causals estructurals}
\perque{Afegeix mecanismes: ja no només descrivim associacions, sinó com es genera cada variable.}

Un SCM és una tupla amb variables exògenes, variables endògenes, funcions
estructurals i una llei per a les exògenes. El canvi conceptual és que ja no
només descrivim distribucions: especifiquem mecanismes que generen les variables.

\diapo{9}{El tabac i el càncer de pulmó}
\perque{Dona un exemple recurrent on associació i efecte causal poden diferir.}

Doll i Hill van intentar establir una relació causal entre fumar i càncer de
pulmó. Fisher va contestar que podia haver-hi un factor genètic que expliqués
totes dues coses.

Aquest exemple ens servirà de fil conductor: si hi ha una causa comuna latent,
el problema ja no és només estadístic. En el graf apareix un \emph{bow}, que és
el patró mínim de confusió latent entre causa i efecte.

\diapo{10}{Operador do()}
\perque{Formalitza la diferència entre observar una variable i forçar-la.}

L'operador \(\operatorname{do}\) formalitza una intervenció. En un SCM, intervenir
vol dir substituir una equació estructural per una assignació constant.

L'exemple de l'escamot d'afusellament ho fa visual: observar que el capità dona
l'ordre no és el mateix que forçar que un tirador decideixi pel seu compte.

\diapo{11}{CBNs Markovianes i semi-Markovianes}
\perque{Distingeix el cas sense confusió latent del cas amb confusió latent.}

Una xarxa bayesiana descriu associacions. Una xarxa bayesiana causal afegeix la
família de distribucions intervingudes. Si és Markoviana, assumim exògenes
independents. Si és semi-Markoviana, podem representar confusió latent amb
arestes bidirigides.

\diapo{12}{El tabac i el càncer de pulmó}
\perque{Mostra per què la porta enrere no resol el cas amb confusió latent.}

Aquí comparem tres lectures: causalitat directa, confusió observada i confusió
latent. Si el confusor fos observat, podríem ajustar-hi. Però en el nostre cas
el factor genètic no és observat.

Per tant, no podem aplicar directament la porta enrere. Necessitem una altra via.

\diapo{13}{Criteri de la porta endavant}
\perque{Ofereix una via alternativa quan hi ha confusió latent entre causa i efecte.}

La porta endavant aprofita un mediador observat \(Z\). La idea és identificar
l'efecte de \(X\) sobre \(Z\), i després l'efecte de \(Z\) sobre \(Y\), ajustant
adequadament.

És més exigent que la porta enrere, però és útil justament quan hi ha confusió
latent entre causa i efecte.

\diapo{14}{El tabac i el càncer de pulmó}
\perque{Fa concreta la porta endavant amb el quitrà com a mediador.}

En l'exemple, el quitrà \(Q\) fa de mediador entre fumar \(T\) i el càncer
\(C\). Si es compleixen les condicions del criteri, podem identificar
\(p(c\mid\operatorname{do}(t))\).

La idea és important: no eliminem la confusió latent, però trobem una manera de
reescriure la pregunta causal amb quantitats observables.

\diapo{15}{do-càlcul}
\perque{Justifica formalment les manipulacions amb intervencions.}

El do-càlcul són tres regles que permeten afegir, eliminar o intercanviar
observacions i intervencions quan les separacions gràfiques adequades es
compleixen.

No cal memoritzar-les ara. El punt és que generalitzen els criteris de porta
enrere i porta endavant.

\diapo{16}{El tabac i el càncer de pulmó}
\perque{Mostra un cas d'èxit: l'efecte causal acaba escrit amb probabilitats observables.}

Aquí es veu la demostració amb do-càlcul de la fórmula de porta endavant.
Primer identifiquem l'efecte de \(T\) sobre \(Q\), i després el de \(Q\) sobre
\(C\).

El resultat final és el que volíem: una quantitat causal expressada només amb
probabilitats observacionals.

\diapo{17}{SUTVA}
\perque{Obre el llenguatge de resultats potencials amb la hipòtesi que fa clara la notació.}

Ara canviem de llenguatge. En el marc de Neyman--Rubin parlem de resultats
potencials: què li passaria a cada unitat sota cada assignació de tractament.

SUTVA diu dues coses: no hi ha interferència entre unitats i el tractament està
ben definit, sense múltiples versions amagades.

\diapo{18}{El problema}
\perque{Presenta el problema fonamental: només observem un resultat potencial per unitat.}

El problema fonamental és que mai observem simultàniament \(Y(1)\) i \(Y(0)\) per
a la mateixa unitat. Si algú rep tractament, observem \(Y(1)\); si no, observem
\(Y(0)\).

Per tant, sempre hi ha resultats potencials no observats. La inferència causal
ha de conviure amb aquest buit.

\diapo{19}{Pearl i Neyman--Rubin}
\perque{Connecta intervencions en SCMs amb resultats potencials.}

Aquí faig explícit el pont. En un SCM, el resultat potencial
\(Y_{\mathbf x}\) és el valor que tindria \(Y\) en el submodel on hem intervingut
\(\mathbf X=\mathbf x\).

Això connecta directament \(P(Y=y\mid\operatorname{do}(\mathbf x))\) amb
\(P(Y_{\mathbf x}=y)\).

\diapo{20}{Pearl i Neyman--Rubin}
\perque{Tanca el pont formal i mostra que els dos llenguatges arriben al mateix estimand.}

Aquest resultat diu que les propietats contrafactuals dels SCMs són suficients
per treballar coherentment amb sistemes causalment ordenats.

A la pràctica, el missatge és que Pearl i Neyman--Rubin no són mons separats:
els SCMs aporten mecanismes i graf; els resultats potencials aporten una lectura
contrafactual i estadística.

\diapo{21}{Berkeley: la lectura agregada}
\perque{Introdueix el cas pràctic amb la lectura aparent de discriminació.}

En el cas Berkeley, l'estudi original tenia \(12763\) sol·licituds a \(101\)
unitats de postgrau. Aquí treballo amb \(4526\) sol·licituds, sis departaments,
admissió, gènere i departament.

La lectura agregada sembla clara: la taxa d'admissió observada afavoreix els
homes.

\diapo{22}{Estratificar en valors absoluts}
\perque{Mostra la distribució real de sol·licituds i admissions per departament.}

Ara estratifiquem per departament i gènere, però en recomptes absoluts.

A l'esquerra veiem que els homes es concentren sobretot als departaments A i B,
mentre que les dones tenen més pes a C, D, E i F. A la dreta veiem admissions
absolutes: aquesta gràfica barreja volum de sol·licituds i selectivitat del
departament.

\diapo{23}{Estratificar en relatius canvia la història}
\perque{Mostra la paradoxa de Simpson: l'agregació pot canviar la lectura.}

Ara miro la mateixa informació en termes relatius. A l'esquerra hi ha la
proporció de sol·licituds dins de cada gènere, i a la dreta les taxes
d'admissió per departament i gènere.

Aquí és on la lectura agregada canvia: el departament pesa molt i no podem
interpretar el total sense mirar l'estratificació.

\diapo{24}{Què passa a nivell associatiu?}
\perque{Recorda que els tests descriuen associació, no causalitat.}

A nivell \(\mathcal L_1\), fem un test d'independència \(\chi^2\) entre gènere i
admissió. Amb \(\alpha=0.05\), a l'agregat rebutgem clarament la independència.

Però departament per departament no passa sempre el mateix: en A rebutgem, en F
no. Això és associació, no causalitat, però ja ens diu que l'agregació amagava
estructura.

\diapo{25}{Quina és la pregunta causal?}
\perque{Mostra que canviar la consulta canvia el significat de la resposta.}

Ara ve el punt causal. No és el mateix intervenir el gènere, observar el
departament, intervenir també el departament, o preguntar què passaria si tothom
sol·licités un departament concret.

I tampoc és el mateix assumir una CBN Markoviana que una semi-Markoviana amb
confusió latent entre departament i admissió.

\diapo{26}{Mateixes dades, diferent identificabilitat}
\perque{Resumeix la lliçó aplicada: la identificabilitat depèn del model causal.}

Amb el model Markovià, totes les consultes de la taula es poden reescriure amb
quantitats observables.

Però si acceptem confusió latent entre departament i admissió, les consultes que
intervenen sobre \(D\) deixen de ser identificables. Les dades són les mateixes;
el que canvia és el model causal.

\diapo{27}{Valors identificats segons el model}
\perque{Concreta amb números quines consultes queden identificades i quines no.}

Aquesta taula posa números a la diapositiva anterior. En el model Markovià, els
valors es poden calcular directament de les dades observades.

En canvi, en el model semi-Markovià amb confusió latent, les consultes que
intervenen sobre el departament apareixen com a no identificables. Això és el
missatge central: no n'hi ha prou amb tenir dades.

\diapo{28}{Conclusions}
\perque{Tanca el relat en una frase central.}

La conclusió és que les dades no contenen causalitat per si soles: contenen
associacions. La causalitat apareix quan hi afegim hipòtesis, coneixement extern
i estructura causal.

Identificar vol dir saber quan una pregunta causal es pot reescriure amb
quantitats observables. Berkeley mostra exactament això: tot depèn de la
pregunta i del model causal assumit.

\diapo{29}{Text final}

Decidir què es mesura i, sobretot, què s'omet; triar els termes adequats; i
filar un relat determinat pot conduir qui ho llegeix, de forma gairebé
imperceptible, cap a una conclusió concreta.

\end{document}
